\begin{lemma}
\label{lem:ch49-truncation-energy}
Let \(\psi\in\PSH(\omega)\), put
\(c:=\sup_{X^{\An}}\psi=\psi(v_{\triv})\), and assume that
\(c\in\mathbf R\) and, for some \(a\ge0\),
\[
                         c-a\ \le\ \psi\ \le\ c.
\]
Then \(\psi\in\cE^1(\omega)\) and
\[
 c-\frac{n}{n+1}\,a\ \le\ E_\omega(\psi)\ \le\ c .
\]
More precisely, writing
\(\mu_j:=V^{-1}(\omega+\ddc\psi)^j\wedge\omega^{\,n-j}\) for the
probability mixed measures of \Cref{thm:bj-mixed-ma}, one has
\(\mu_0=\delta_{v_{\triv}}\) and
\[
 (n+1)\bigl(E_\omega(\psi)-c\bigr)
 =\sum_{j=1}^{n}\int_{X^{\An}}(\psi-c)\,\rd\mu_j .
\]
In particular, one may take
\(a=\operatorname{osc}_{X^{\An}}\psi=c-\inf_{X^{\An}}\psi\).  When
\(\sup\psi=0\), this gives
\(-\frac{n}{n+1}a\le E_\omega(\psi)\le0\).
\end{lemma}

\begin{proof}
We first make explicit the precise results from
\Cref{chap:bj-energy,chap:bj-finite-energy} that enter the argument.  The
normalized energy is monotone and satisfies
\[
 E_\omega(t)=t,
 \qquad
 E_\omega(\eta+t)=E_\omega(\eta)+t,
 \qquad
 \eta\le\chi\Longrightarrow E_\omega(\eta)\le E_\omega(\chi).
\tag{49.6a}\label{eq:ch49-truncation-input-energy}
\]
The extended \((n+1)\)-fold pairing is continuous along coordinatewise
decreasing nets.  If \(\eta\in\cE^1(\omega)\), its normalized mixed
measures satisfy, for \(0\le j\le n\),
\[
 \begin{aligned}
 \nu_j(\eta)&:=V^{-1}(\omega+\ddc\eta)^j\wedge\omega^{\,n-j}\ge0,
 &\nu_j(\eta)(X^{\An})&=1,\\
 \int f\,\rd\nu_j(\eta)
 &=V^{-1}(0,f)(\omega,\eta)^j(\omega,0)^{n-j}
 &&(f\in\cE^1(\omega)).
 \end{aligned}
\tag{49.6b}\label{eq:ch49-truncation-input-mixed}
\]
These are, respectively,
\Cref{prop:ch46-energy-usc,thm:bj-extended-energy-pairing,thm:bj-mixed-ma};
the formula in \eqref{eq:ch49-truncation-input-mixed} also fixes every
normalizing factor needed below.

The bounds on \(\psi\), together with
\eqref{eq:ch49-truncation-input-energy}, give
\[
 c-a=E_\omega(c-a)\le E_\omega(\psi)\le E_\omega(c)=c.
\]
Thus \(E_\omega(\psi)\) is finite and
\(\psi\in\cE^1(\omega)\), without any appeal to an unrecorded
integrability criterion.  We also recall the short proof of the maximum
principle in the form used here.  Choose a decreasing Fubini--Study
regularization \(\psi_\alpha\downarrow\psi\) in the fixed ample class
\(\omega\).  Each \(\psi_\alpha\) attains its maximum at
\(v_{\triv}\): in a weighted-section presentation every nonzero section
satisfies \(\log|s|\le0\) on \(X^{\An}\) and
\(\log|s|(v_{\triv})=0\), so the weighted maximum is everywhere at most
its value at \(v_{\triv}\).  Consequently
\[
 \sup\psi\le\inf_\alpha\sup\psi_\alpha
 =\inf_\alpha\psi_\alpha(v_{\triv})
 =\psi(v_{\triv})\le\sup\psi,
\]
and hence \(c=\psi(v_{\triv})\).

For completeness, we now pass the PL telescoping identity to \(\psi\),
rather than merely citing the resulting mixed-energy formula.  At every
Fubini--Study level, multilinearity gives
\[
 (\omega,\psi_\alpha)^{n+1}-(\omega,0)^{n+1}
 =\sum_{j=0}^n
   (0,\psi_\alpha)(\omega,\psi_\alpha)^j(\omega,0)^{n-j}.
\tag{49.6c}\label{eq:ch49-truncation-telescope-pl}
\]
For each \(j\), the summand on the right is the difference
\[
 (\omega,\psi_\alpha)^{j+1}(\omega,0)^{n-j}
 -(\omega,\psi_\alpha)^j(\omega,0)^{n-j+1}.
\]
Thus coordinatewise decreasing continuity of the extended pairing passes
\eqref{eq:ch49-truncation-telescope-pl} to the limit without subtracting
two infinite quantities: all limiting pairings are finite because
\(\psi\) is bounded.  Moreover \((\omega,0)^{n+1}=0\), since on the
product compactification it is the intersection of \(n+1\) pullbacks from
the \(n\)-fold \(X\).  Dividing the limiting
identity by \(V\), using
\((n+1)E_\omega(\psi)=V^{-1}(\omega,\psi)^{n+1}\), and then applying
\eqref{eq:ch49-truncation-input-mixed} with \(f=\psi\), gives
\[
 (n+1)E_\omega(\psi)
 =\sum_{j=0}^n\int_{X^{\An}}\psi\,\rd\mu_j .
\tag{49.6d}\label{eq:ch49-truncation-mixed-energy}
\]

It remains to identify the zeroth term and estimate the others.  On the
trivial test configuration, the central fibre is one reduced component,
its associated point is \(v_{\triv}\), and its intersection coefficient
is \((\omega^n)=V\).  Hence the unnormalized measure \(\omega^n\) is
\(V\delta_{v_{\triv}}\), so
\[
                         \mu_0=V^{-1}\omega^n=\delta_{v_{\triv}}.
\]
The mass-one assertion in \eqref{eq:ch49-truncation-input-mixed} now gives
\(\int c\,\rd\mu_j=c\) for every \(j\), while the preceding maximum
principle gives \(\int\psi\,\rd\mu_0=c\).  Subtracting \((n+1)c\) from
\eqref{eq:ch49-truncation-mixed-energy} therefore yields
\[
 (n+1)\bigl(E_\omega(\psi)-c\bigr)
 =\sum_{j=1}^n\int_{X^{\An}}(\psi-c)\,\rd\mu_j.
\]
Every integral on the right lies in \([-a,0]\), since
\(-a\le\psi-c\le0\) and \(\mu_j\) is a probability measure.  Division
by \(n+1\) proves the asserted bounds.  Taking the actual oscillation gives
the first specialization.  Finally, \(\ddc(\psi+t)=\ddc\psi\) and
\eqref{eq:ch49-truncation-input-energy} is translation covariant, so
adding \(t\) changes both \(c\) and \(E_\omega(\psi)\) by \(t\) and
leaves \(a\) and all \(\mu_j\) unchanged; this gives the supremum-zero
form as well.
\end{proof}
